\begingroup
\makeatletter\tagsleft@true\let\veqno\@@leqno\makeatother
\clearpage\NoSourceHead\SourceLeaf{345}
\SourceGroup{III}
\begin{center}\large\bfseries EINIGES ÜBER DIE UNENDLICHE REIHE\end{center}
\[\begin{aligned}
&1+\frac{\alpha{.}\beta}{1{.}\gamma}x+\frac{\alpha{.}\alpha+1{.}\beta{.}\beta+1}{1{.}2{.}\gamma{.}\gamma+1}xx\\[.5em]
&\quad+\frac{\alpha{.}\alpha+1{.}\alpha+2{.}\beta{.}\beta+1{.}\beta+2}{1{.}2{.}3{.}\gamma{.}\gamma+1{.}\gamma+2}x^3{.}{.}{.}{.}
\end{aligned}\]
\HistoricCitation{[Aus Handbuch 19, Be, Kleine Aufsätze aus verschiedenen Theilen der Mathematik,\\Angefangen im May 1809.]}
\NumberedItem{1}
\noindent{\small[S. 36]}\par
\par\indent Wir bezeichnen den Werth dieser Reihe, welche stets convergirt, wenn $x$ kleiner ist als 1, durch das Zeichen $F(\alpha,\beta,\gamma)$, so wie die davon derivirte Function oder $\frac{dF(\alpha,\beta,\gamma)}{dx}$ durch $F'(\alpha,\beta,\gamma)$. Wir haben also sofort
\begin{equation}
F'(\alpha,\beta,\gamma)=\frac{\alpha\beta}{\gamma}F(\alpha+1,\beta+1,\gamma+1).\tag*{(1)}
\end{equation}
\par\indent Die Coefficienten von $x^r$ verhalten sich bei den drei Functionen
\[F(\alpha,\beta,\gamma),\quad xF'(\alpha,\beta,\gamma),\quad F(\alpha+1,\beta,\gamma)\]
\noindent offenbar wie
\[1,\qquad r,\qquad \frac{\alpha+r}{\alpha}\mathpunct{\raisebox{.45ex}{.}}\]
\noindent Wir schliessen hieraus
\begin{equation}
F(\alpha+1,\beta,\gamma)=F(\alpha,\beta,\gamma)+\frac{x}{\alpha}F'(\alpha,\beta,\gamma)\tag*{(2)}
\end{equation}
\noindent und auf ähnliche Art
\begin{equation}
F(\alpha,\beta+1,\gamma)=F(\alpha,\beta,\gamma)+\frac{x}{\beta}F'(\alpha,\beta,\gamma).\tag*{(3)}
\end{equation}
\clearpage\SourceLeaf{346}\ArithmeticHead{339}{R}{{\scriptsize ZUR THEORIE DER REIHE $F(\alpha,\beta,\gamma,x)$.\quad III.}}
\noindent Hingegen werden sich die Coefficienten derselben Potenz von $x$ in $F'(\alpha,\beta,\gamma)$ und $F'(\alpha+1,\beta,\gamma)$ noch wie
\[\frac{\alpha+r{.}\beta+r}{\gamma+r}\quad\text{und}\quad\frac{(\alpha+r+1)(\alpha+r)(\beta+r)}{\alpha(\gamma+r)}\]
\noindent und der von $xF'(\alpha+1,\beta,\gamma)$ wie $\frac{(\alpha+r)r}{\alpha}$ verhalten. Der Coefficient in
\[(1-x)F'(\alpha+1,\beta,\gamma)\]
\noindent verhält sich also wie
\[\frac{(\alpha+r)(\alpha\beta+\beta+(\alpha+\beta-\gamma+1)r)}{\alpha(\gamma+r)},\]
\noindent folglich der in
\[(\alpha+1-\gamma)F'(\alpha,\beta,\gamma)-\alpha(1-x)F'(\alpha+1,\beta,\gamma)\]
\noindent wie $-\beta(\alpha+r)$. Wir haben also
\par\smallskip\noindent{\small[S. 37]}\par
\begin{equation}
F'(\alpha+1,\beta,\gamma)=\frac{\beta}{1-x}F(\alpha,\beta,\gamma)+\frac{\alpha-\gamma+1+\beta x}{\alpha(1-x)}F'(\alpha,\beta,\gamma)\tag*{(4)}
\end{equation}
\noindent und eben so ist offenbar
\begin{equation}
F'(\alpha,\beta+1,\gamma)=\frac{\alpha}{1-x}F(\alpha,\beta,\gamma)+\frac{\beta-\gamma+1+\alpha x}{\beta(1-x)}F'(\alpha,\beta,\gamma).\tag*{(5)}
\end{equation}
\noindent Aus der Verbindung von (2) und (4) schliessen wir ferner
\begin{equation}
F(\alpha-1,\beta,\gamma)=\frac{\alpha-\gamma+\beta x}{\alpha-\gamma}F(\alpha,\beta,\gamma)-\frac{x-xx}{\alpha-\gamma}F'(\alpha,\beta,\gamma)\tag*{(6)}
\end{equation}
\noindent und eben so
\begin{equation}
F(\alpha,\beta-1,\gamma)=\frac{\beta-\gamma+\alpha x}{\beta-\gamma}F(\alpha,\beta,\gamma)-\frac{x-xx}{\beta-\gamma}F'(\alpha,\beta,\gamma),\tag*{(7)}
\end{equation}
\begin{equation}
F'(\alpha-1,\beta,\gamma)=-\frac{(\alpha-1)\beta}{\alpha-\gamma}F(\alpha,\beta,\gamma)+\frac{(\alpha-1)(1-x)}{\alpha-\gamma}F'(\alpha,\beta,\gamma),\tag*{(8)}
\end{equation}
\begin{equation}
F'(\alpha,\beta-1,\gamma)=-\frac{(\beta-1)\alpha}{\beta-\gamma}F(\alpha,\beta,\gamma)+\frac{(\beta-1)(1-x)}{\beta-\gamma}F'(\alpha,\beta,\gamma).\tag*{(9)}
\end{equation}
\par\indent Hieraus ist klar, dass man aus $F(\alpha,\beta,\gamma)$ und $F'(\alpha,\beta,\gamma)$ allgemein
\[F(\alpha+k,\beta+l,\gamma)\quad\text{und}\quad F'(\alpha+k,\beta+l,\gamma)\]
\noindent für jede ganze Werthe von $k$ und $l$, sie mögen positiv oder negativ seyn, rational ableiten kann.
\par\bigskip\noindent\hfill 43\textsuperscript{*}
\clearpage
\endgroup
